\scanpage{76}\mspage{63}
% archive sheet number: 75
% Paragraph continues from the preceding sheet.
on peut appliquer les résultats précédents,
notamment 3.5., aux flèches $v,w$. \deleted{\unsure{Le}}
\unsure{Lemme|Le lemme} précédent montre bien que $v$
satisfait aux conditions de 3.5, donc
le foncteur $v_!^W$ construit ds 3.3.1. est
adjoint à gauche de $v_W^*$. \inserted{Considérons} \deleted{\unsure{D}}d'autre part,
le foncteur $w_!^W$ construit \unsure{ds|dans} (3.3.1), je
dis que la $\unsure{=}$ conclusion s'applique.

Cela provient du fait que $w$ est une
immersion fermée. En effet, dans ce
cas il est immédiat que \deleted{$w_!:\unsure{u}$}
\[
w_!:M(I^{\unsure{o|\circ}})\unsure{\text{---}|\to} M(\deleted{\unsure{I}}J^o)\qquad(\text{où }J=C(u))
\]
est compatible avec les localiseurs,
puisque \unsure{plus|fibre} \unsure{à|a}
\[
\left\{\begin{aligned}
 w_!(X)(j)&\simeq X(j) &&\text{si }\deleted{\unsure{i}}j\in I\\
 w_!(X)(j)&\simeq\varnothing_M &&\text{si }j\in J\setminus I\,,
\end{aligned}\right.
\]
par le calcul des fibres.
\[
w_!(X)(j)\simeq\underrightarrow{\lim}_{i\in I/j}X(i)\,.
\]

Donc \unsure{le|la} \deleted{calcul}\inserted{définition} de $w_!^W$ \unsure{donné} dans
3.3.1. revient en fait à passer directement
aux cat. localisées pour $w_!$, sans avoir
à se restreindre à $M(I^o)_c$, et le
lemme des foncteurs adjoints nous dit que
